%%% FORMATO E IDIOMA %%%
    \documentclass[letterpaper,DIV=15,12pt]{scrartcl}
    \usepackage[spanish,mexico,shorthands=off,es-noindentfirst]{babel}
    \usepackage{scrlayer-scrpage}
        \clearpairofpagestyles
        \ihead{\footnotesize \textit{Teoría de Conjuntos III}}
        \ohead{\footnotesize \textit{2026-2}}
        \cfoot{\normalfont\thepage}
        \addtokomafont{title}{\bfseries \rmfamily}
        \setlength{\headsep}{5pt}
        \newpairofpagestyles{beginstyle}{
            \clearpairofpagestyles
            \KOMAoptions{headsepline=false}
            \cfoot{\footnotesize \pagemark}
            \cfoot{\normalfont\thepage}
        }
        \addtokomafont{section}{\raggedleft\rmfamily}
        \addtokomafont{subsection}{\raggedleft\rmfamily}
        \addtokomafont{subsubsection}{\raggedleft\rmfamily}
        \setlength{\headsep}{8pt}
        \setlength{\footskip}{20pt}
        \setlength{\textheight}{600pt}
%%% PAQUETES %%%
    \usepackage[style=mexican]{csquotes}
    \usepackage{ifthen}
    \usepackage[shortlabels]{enumitem}
        \setenumerate[1]{label=\MakeLowercase{\roman*}), ref=\roman*}
        \setenumerate[2]{label=\MakeLowercase{\alph*}), ref=\alph*}
    \usepackage{mathtools}
    \usepackage{amsthm}
    \usepackage[T1]{fontenc}
    \usepackage{stix2}
	\renewcommand{\qedsymbol}{$\blacksquare$}
	\addto\captionsspanish{\renewcommand{\proofname}{\normalfont\bfseries Demostración}}	
	\newenvironment{solucion}{\renewcommand{\qedsymbol}{$\boxtimes$}\begin{proof}[\normalfont\bfseries Solución]}{\end{proof}}

        \newtheoremstyle{manual}{3pt}{3pt}{\normalfont}{}{\bfseries}{.}{.5em}{#1 #3}
		\theoremstyle{manual}
		\newtheorem{proposicion}{Proposición}
        \newtheorem{definicion}{Definición}
        \newtheorem{teorema}{Teorema}
        \newtheorem{lema}{Lema}
        \newtheorem{corolario}{Corolario}
        \newtheorem*{ejercicio*}{Ejercicio}        
%%% C O M A N D O S %%%
    \renewcommand{\emptyset}{\varnothing}
    \NewDocumentCommand{\mb}{O{}}{\cramped{V^{(B)}_{#1}}}      %V^(B)
    \NewDocumentCommand{\mg}{O{}}{\cramped{V^{(\Gamma)}_{#1}}} %V^(Γ)
    \DeclarePairedDelimiter{\vb}{\lBrack}{\rBrack}             %|| - ||
    \DeclareMathOperator{\dom}{dom}                            %dominio
    \DeclareMathOperator{\cod}{cod}                            %codominio
    \DeclareMathOperator{\im}{im}                              %imagen 
%conjuntos
\newcommand{\st}{}                                              
\newcommand{\SetSymbol}[1][]{%
\nonscript\:#1\vert
\allowbreak
\nonscript\:
\mathopen{}}
\DeclarePairedDelimiterX\set[1]\{\}{%
\renewcommand\st{\SetSymbol[\delimsize]}
#1
}
%% D O C U M E N T O %%

\begin{document}

\thispagestyle{plain}

    \begin{center}
        {\fontsize{30}{60}\rmfamily \textbf{Modelo booleano-valuado}} \\ \vspace{.2cm} Teoría de Conjuntos III, 2026-2
    \end{center}
    \begin{flushright}
        \footnotesize \hfill Profesor: Luis Jesús Turcio Cuevas.\\
        \hfill Ayudante: Hugo Víctor García Martínez.
    \end{flushright}

\section{Validez de los axiomas}
Recordemos que tenemos un álgebra de Boole completa \(B\) y la construcción de
una estructura booleano-valuada dada por la siguiente jerarquía acumulativa:
\begin{align*}
  \mb[0] &=\emptyset,\\
  \mb[\alpha+1] &=\set{f\colon X\to B\st X\subseteq\mb[\alpha]},\\
  \mb[\gamma] &=\bigcup_{\alpha<\gamma}\mb[\alpha],\\
  \mb &=\smashoperator{\bigcup_{\alpha\in\text{OR}}}\mb[\alpha].
\end{align*}
Además, definimos el valor booleano de las fórmulas atómicas de manera
simultánea y por recursión de la siguiente manera:
\begin{align*}
  \vb{x\in y} &= \smashoperator{\sum_{s\in\dom(y)}}y(s)\cdot\vb{s=x},\\
  \vb{x\subseteq y} 
  &= \smashoperator{\prod_{t\in\dom(x)}}x(t)\Rightarrow\vb{t\in y},\\
  \vb{x=y} &= \vb{x\subseteq y}\cdot\vb{y\subseteq x}.
\end{align*}
Con esta definición de valor booleano se satisface, en particular, la
sustitución de iguales por iguales.
Por último recordemos que \(\mb\models\varphi\) si y sólo si \(\vb{\varphi}=1\).


\subsection{Extensionalidad}
Sean \(x,y\in\mb\) y veamos que
\begin{equation}
  \label{eq:ext}
  \vb{\forall w(w\in x\to w\in y)\to x\subseteq y}=1.
\end{equation}
Con la otra implicación del axioma obtendríamos la otra contención. Como la
demostración del enunciado anterior es análoga a lo que veremos, la omitiremos.

Por la equivalencia del orden con la implicación booleana, para ver que se
satisface la ecuación~\ref{eq:ext} es suficiente mostrar que
\begin{equation}
  \vb{\forall w(w\in x\to w\in y)}\leq\vb{x\subseteq y}.
\end{equation}

Consideramos la siguiente cadena de desigualdades
\begin{align*}
  \vb{\forall w(w\in x\to w\in y)} 
  &= \smashoperator{\prod_{a\in\mb}}\vb{a\in x\to a\in y}
  && \text{definición de \(\vb{\forall\dots}\)}\\
  &\leq \smashoperator{\prod_{t\in\dom(x)}}\vb{t\in x\to t\in y}
  && \text{propiedades de ínfs}\\
  &= \smashoperator{\prod_{t\in\dom(x)}}\vb{t\in x}\Rightarrow \vb{t\in y}\\
  &\leq \smashoperator{\prod_{t\in\dom(x)}}x(t)\Rightarrow\vb{t\in y}
  && x(t)\leq\vb{t\in x}\\
  &= \vb{x\subseteq y}.
\end{align*}


\subsection{Separación}
Sean \(x\in\mb\) y \(\varphi(v)\) una fórmula. Sea \(z\in\mb\) que satisface
\(\dom(z)=\dom(x)\) y \(z(t)=x(t)\cdot\vb{\varphi(t)}\). Veamos que se satisface
\begin{equation*}
  \vb{\forall w(w\in z \leftrightarrow w\in x\land\varphi(w))}=1
\end{equation*}
Como el valor booleano de un cuantificador universal es un ínfimo y queremos que
este sea \(1\), entonces cada uno de los \textit{ínfimandos} debe ser \(1\).
Así, tenemos que mostrar que para cualquier \(a\in\mb\) se satisface
\begin{equation*}
  \vb{a\in z \leftrightarrow a\in x\land\varphi(a)}=1.
\end{equation*}
Por la definición del bicondicional y la equivalencia del orden con implicación
booleana, lo anterior se traduce a demostrar
\begin{equation*}
  \vb{a\in z}=\vb{a\in x\land\varphi(a)}.
\end{equation*}
Consideremos la siguiente cadena de igualdades
\begin{align*}
  \vb{a\in z} &= \smashoperator{\sum_{t\in\dom(z)}}z(t)\cdot\vb{t=a}\\
  &= \smashoperator{\sum_{t\in\dom(x)}}x(t)\cdot\vb{\varphi(t)}\cdot\vb{t=a}
  && \text{definición de \(z\)}\\
  &= \smashoperator{\sum_{t\in\dom(x)}}x(t)\cdot\vb{\varphi(a)}\cdot\vb{t=a}
  && \text{sustitución de = por =}\\
  &= \left(\smashoperator[r]{\sum_{t\in\dom(x)}}x(t)
  \cdot\vb{t=a}\right)\cdot\vb{\varphi(a)}\\
  &= \vb{a\in x}\cdot\vb{\varphi(a)}.
\end{align*}

Como ya tenemos separación, ahora es posible verificar la versión débil de los
axiomas. Por ejemplo, para el axioma del par, si \(x,y\in\mb\) entonces es
suficiente dar la existencia de \(z\in\mb\) tal que las fórmulas \(x\in z\) y
\(y\in z\) se satisfagan en \(\mb\), es decir, un supraconjunto del par de
verdad. 


\subsection{Par}
Sean \(x,y\in\mb\). Definimos \(z\in\mb\) como sigue: \(\dom(z)=\set{a,b}\) y
\(z(a)=1=z(b)\). Aunque no es difícil ver que con esta definición \(z\)
satisface la versión original del axioma del par, sólo veremos que 
\begin{equation*}
  \vb{x\in z}=\vb{y\in z}=1.
\end{equation*}
Además, la demostración con \(y\) es análoga a la de \(x\) así que sólo haremos
una.
\begin{align*}
  \vb{x\in z} &= (z(x)\cdot\vb{x=x})+(z(y)\cdot\vb{x=y})\\
  &= 1 + (z(y)\cdot\vb{x=y})\\
  &=1.
\end{align*}


\subsection{Unión}
Sea \(x\in\mb\) y proponemos \(z\in\mb\) tal que
\begin{align*}
  \dom(z) &= \smashoperator{\bigcup_{t\in\dom(x)}}\dom(t)\\
  z(r) &= 1.
\end{align*}
Veamos que \(\vb{\forall y(y\in x\to \forall w(w\in y\to w\in z))}=1\). La
fórmula dentro del valor booleano anterior se puede abreviar usando
cuantificadores acotados. Con la abreviatura mencionada lo que tenemos que
mostrar es:
\begin{equation*}
  \vb{\forall y\in x\forall w\in y(x\in z)}=1.
\end{equation*}
Como la anterior es una fórmula con cuantificadores acotados, entonces la
ecuación anterior se traduce a
\begin{equation*}
  \smashoperator{\prod_{t\in\dom(x)}} x(t)\Rightarrow\vb{\forall w\in t(w\in z)}=1,
\end{equation*}
es decir, dada \(t\in\dom(x)\) veamos que 
\(x(t)\leq\vb{\forall w\in t(w\in z)}\). De nuevo, el valor booleano del lado
derecho de la implicación es un cuantificador acotado que sabemos es igual a
\(\prod_{r\in\dom(t)}t(r)\Rightarrow\vb{r\in z}\). Por las propiedades de
ínfimo, todo lo anterior se traduce a demostrar que si \(t\in\dom(x)\) y
\(r\in\dom(t)\), entonces
\begin{equation*}
  x(t)\leq (t(r)\Rightarrow\vb{r\in z}).
\end{equation*}
Por la adjunción entre ínfimo e implicación veremos que 
\(x(t)\cdot t(r)\leq\vb{r\in z}\). Para esto notamos que \(r\in\dom(z)\) por lo
que \(\vb{r\in z}\geq z(r)=1\) y así lo que queremos demostrar es inmediato.


\subsection{Infinito}
Como \enquote{\(x\) es inductivo} es una fórmula \(\Delta_{0}\), entonces
\begin{equation*}
  \vb{\check{\omega}\text{ es inductivo}}=1
\end{equation*}
mostrando así que hay un \enquote{conjunto} inductivo en \(\mb\).


\subsection{Buena fundación}
Sea \(x\in\mb\) y veamos que si \(x\ne\emptyset\) entonces tiene un
\(\in\)-minimal, es decir,
\begin{equation*}
  \vb{\exists u(u\in x)\to\exists y\in x\forall z\in y(z\notin x)}=1.
\end{equation*}
Supongamos que no es cierto, es decir,
\begin{equation*}
  \vb{\exists u(u\in x)\land\forall y\in x\exists z\in y(z\in x)}=b\ne 0.
\end{equation*}

Como estamos en una jerarquía acumulativa, entonces podemos tomar \(a\in\mb\) de
rango mínimo tal que \(\vb{a\in x}\cdot b\ne 0\). Como \(b\) es un ínfimo,
entonces es menor o igual a sus \enquote{ínfimandos} y hacer ínfimo con \(b\)
preserva el orden. Con lo anterior tenemos:
\begin{equation*}
  \vb{a\in x}\cdot b\leq
  \vb{a\in x}\cdot\vb{\forall y\in x\exists z\in y(z\in x)}\cdot b.
\end{equation*}
Notemos que las fórmulas que están escritas explícitamente en el lado derecho de
la desigualdad anterior son \(\Delta_{0}\). Además, la conjunción de las dos implica a
la fórmula \(\exists z\in a(z\in x)\), que también es \(\Delta_{0}\). Así,
tenemos una implicación de fórmulas \(\Delta_{0}\) válida en \(V\), por lo que
el valor booleano de la implicación debe ser \(1\) en \(\mb\), o bien tenemos la
siguiente desigualdad
\begin{equation*}
  \vb{a\in x}\cdot\vb{\forall y\in x\exists z\in y(z\in x)}\cdot b \leq
  \vb{\exists z\in a(z\in x)}\cdot b.
\end{equation*}
De la condición que satisface \(a\), la distributividad general y el valor de un
cuantificador existencial acotado tenemos
\begin{equation*}
  \smashoperator{\sum_{t\in\dom(a)}}(a(t)\cdot\vb{t\in x}\cdot b)\ne 0.
\end{equation*}
Por lo que debe existir \(t\in\dom(a)\) tal que 
\(a(t)\cdot\vb{t\in x}\cdot b\ne 0\). Con lo cual tenemos que el rango de \(t\)
es estrictamente menor que el de \(a\) y \(\vb{t\in x}\cdot b\ne 0\), lo cual
contradice nuestra elección de \(a\).


\subsection{Reemplazo}
Sean \(x\in\mb\) y \(\varphi\) una fórmula con dos variables libres que se
comporta de manera funcional, al menos en \(x\), es decir, satisface
\begin{equation*}
  \forall u(u\in x\to \exists! v\, \varphi(u,v)).
\end{equation*}
Veamos que existe \(z\in\mb\) que cumple
\begin{equation}
  \label{eq:re1}
  \vb{\forall u\in x(\exists v\,\varphi(u,v)\to\exists v\in z\,\varphi(u,v))}
  =1.
\end{equation}

Definimos \(z\in\mb\) mediante \(\dom(z)=\bigcup\set{S_t\st\in\dom(x)}\) y
\(z(s)=1\), donde \(S_t\) es un conjunto (por axioma de elección) que satisface
\begin{equation*}
  \sum_{y\in\mb}\vb{\varphi(t,y)}=\sum_{y\in S_{t}}\vb{\varphi(t,y)}.
\end{equation*}
Por las fórmulas que dicen cómo evaluar el valor booleano de cuantificadores
acotados, para ver que se satisface la ecuación~\eqref{eq:re1} debemos mostrar
\begin{equation*}
  x(t)\leq\vb{\exists v\,\varphi(t,v)\to\exists v\in z\,\varphi(t,v)},
\end{equation*}
para toda \(t\in\dom(x)\). Esto lo haremos viendo que el lado derecho es igual a
\(1\), es decir, veremos que
\begin{equation*}
  \vb{\exists v\,\varphi(t,v)}\leq\vb{\exists v\in z\,\varphi(t,v)}.
\end{equation*}
Procedemos mediante la siguiente cadena de desigualdades
\begin{align*}
  \vb{\exists v\,\varphi(t,v)} &=
  \smashoperator{\sum_{a\in\mb}}\vb{\varphi(t,a)}\\
  &= \smashoperator{\sum_{a\in S_t}}\vb{\varphi(t,a)}
  && \text{definición de \(S_t\)}\\
  &\leq \smashoperator{\sum_{a\in\dom(z)}}\vb{\varphi(t,a)}
  && \text{definición de \(\dom(z)\)}\\
  &= \smashoperator{\sum_{a\in\dom(z)}}z(a)\cdot\vb{\varphi(t,a)}
  && \text{definición de \(z(s)\)}\\
  &= \vb{\exists v\in z\,\varphi(t,v)}.
\end{align*}


\subsection{Elección}
Recordemos que un conjunto booleano \(x\in\mb\) lleva la información de sus
elementos en \(\dom(x)\) y la regla de correspondencia \(x(t)\) nos dice que
tanto \(t\) es elemento de \(x\). Así, para dar una función \(f\) en \(\mb\)
tenemos que hacer que el dominio de \(f\) sea un conjunto de pares ordenados
\((a,b)\) según \(\mb\) y la regla de correspondencia \(f(a,b)\) no dirá que
tanto \(a\) va a dar a \(b\) bajo la función \(f\). Siguiendo estas ideas, la
imagen de \(f\) en \(\mb\) estará dada por
\begin{equation}
  \label{eq:im}
  \begin{gathered}
    \dom(\im(f))=\set{b\st(a,b)\in\dom(f)},\\
    \im(f)(b)=\smashoperator{\sum_{(a,b)\in\dom(f)}}f(a,b).
  \end{gathered}
\end{equation}

Ahora sí, veamos que se vale elección en \(\mb\).
Primero notemos que la formula \enquote{\(s\) es bien ordenable} es
\(\Sigma_{1}\). Por axioma de elección en \(V\) tenemos que todo conjunto es
bien ordenable y así
\begin{equation*}
  \vb{\check{s}\text{ es bien ordenable}}=1.
\end{equation*}
Sea \(x\in\mb\) y veamos que existe \(s\in V\) y una función (según \(\mb\))
\(f\in\mb\) que cumplen
\begin{equation*}
  \vb{\text{\(f\) es una función sobre \(\check{s}\) y
  \(x\subseteq\im(f)\)}}=1.
\end{equation*}

Tomamos \(s=\dom(x)\in V\) y consideramos la función
\(g\colon\dom(\check{s})\to s\) definida por \(g(\check{y})=y\), para todo
\(y\in s\). Como el dominio y codominio de \(g\) están en \(V\), entonces 
\(g\in V\) y además es claro que \(g\) es una función sobre \(s\)
(suprayectiva). Como la fórmula que dice \enquote{\(g\) es sobre \(s\)} es
\(\Delta_{0}\), entonces
\begin{equation*}
  \vb{\text{\(\check{g}\) es una función sobre \(\check{s}\)}}=1.
\end{equation*}

Para terminar veamos que \(\vb{x\subseteq\im(f)}=1\). Por la definición del
valor booleano de la inclusión, lo anterior se traduce a demostrar que para todo
\(t\in\dom(x)=s\) se tiene que \(x(t)\leq\vb{t\in\im(\check{g})}\). De hecho,
veremos que el lado derecho de la desigualdad anterior es igual a \(1\).
\begin{align*}
  \vb{t\in\im(\check{g})} 
  &= \smashoperator{\sum_{(\check{y},y)\in\dom(\check{g})}}
  \im(\check{g})(y)\cdot\vb{y=t}\\
  &= \smashoperator{\sum_{(\check{y},y)\in\dom(\check{g})}}\vb{y=t}
  &&\text{Por~\eqref{eq:im} y definición de \(\check{g}\)}\\
  &= \smashoperator{\sum_{y\in s}}\vb{y=t}
  &&\text{Por definición de \(\check{g}\) y \(s\)}\\
  &\geq \vb{t=t}\\
  &= 1.
\end{align*}


\subsection{Potencia}
Sea \(x\in\mb\) y veamos que existe \(z\in\mb\) que cumple \(\vb{\forall
y(y\subseteq x\to y\in z)}=1.\) Definimos \(z\) como sigue:
\begin{gather*}
  \dom(z)=\set{u\in\mb\st\dom(u)=\dom(x)}=B^{\dom(x)}\\
  z(u)=\vb{u\subseteq x}.
\end{gather*}
Veamos que \(z\) cumple lo que queremos, es decir, que si \(y\in\mb\) entonces
\(\vb{y\subseteq x}\leq\vb{y\in z}\). Para esto daremos \(y'\in\mb\) tal que
\begin{equation*}
  \vb{y\subseteq x}\leq\vb{y'\in z}\cdot\vb{y=y'}.
\end{equation*}
Definimos \(y'\) con \(\dom(y')=\dom(x)\), es decir, \(y'\in\dom(z)\) y
\(y'(t)=\vb{t\in y}\).

Primero notemos que se satisface \(\vb{y'\subseteq y}=1\), como se muestra a
continuación
\begin{align*}
  \vb{y'\subseteq y} 
  &= \smashoperator{\prod_{t\in\dom(y')}}y'(t)\Rightarrow\vb{t\in y}\\
  &= \smashoperator{\prod_{t\in\dom(x)}}\vb{t\in y}\Rightarrow\vb{t\in y}\\
  &= 1.
\end{align*}

Además de lo anterior también se cumple \(\vb{x\cap y\subseteq y'}=1\). Para
mostrar la identidad anterior escribimos de manera formal la fórmula dentro del
valor booleano, \(\forall w(w\in x\cap y\to w\in y')\). Por lo tanto, es
suficiente demostrar que si \(w\in\mb\) entonces 
\(\vb{w\in x\cap y}\leq\vb{w\in y'}\).

\begin{align*}
  \vb{w\in x\cap y} &= \vb{w\in x}\cdot\vb{w\in y}\\
  &= \smashoperator{\sum_{t\in\dom(x)}}x(t)\cdot\vb{t=w}\cdot\vb{w\in y}
  && \text{definición de \(\vb{w\in x}\)}\\
  &\leq \smashoperator{\sum_{t\in\dom(x)}}\vb{t=w}\cdot\vb{t\in y}\\
  &= \smashoperator{\sum_{t\in\dom(x)}}\vb{t=w}\cdot y'(t)
  && \text{definición de \(y'\)}\\
  &= \vb{w\in y'}.
\end{align*}

Ahora usaremos que la fórmula que dice 
\enquote{\(a\) está contenido en \(b\)} es \(\Delta_{0}\) para concluir:
\begin{align*}
  \vb{y\subseteq x} &= \vb{y\subseteq x}\cdot\vb{x\cap y\subseteq y'}\\
  &= \vb{y\subseteq x\land x\cap y\subseteq y'}\\
  &\leq \vb{y\subseteq y'}\\
  &= \vb{y\subseteq y'}\cdot\vb{y'\subseteq y}\\
  &= \vb{y=y'}.
\end{align*}

Finalmente notamos que se tiene la siguiente desigualdad
\begin{align*}
  \vb{y\subseteq x} &= \vb{\forall w(w\in y\to w\in x)}\\
  &= \smashoperator{\prod_{a\in\mb}}\vb{a\in y}\Rightarrow\vb{a\in x}\\
  &\leq \smashoperator{\prod_{a\in\dom(y')}}\vb{a\in y}\Rightarrow\vb{a\in x}\\
  &= \smashoperator{\prod_{a\in\dom(y')}}y'(a)\Rightarrow\vb{a\in x}\\
  &= \vb{y'\subseteq x}\\
  &= z(y')\\
  &\leq \vb{z'\in y}.
\end{align*}

Con todo lo anterior hemos demostrado
\begin{equation*}
  \vb{y\subseteq x}\leq\vb{y'\in z}\cdot\vb{y=y'}\leq\vb{y\in z}.
\end{equation*}

\end{document}